% Part: computability
% Chapter: computability-theory
% Section: normal-form

\documentclass[../../../include/open-logic-section]{subfiles}

\begin{document}

\olfileid{cmp}{thy}{nfm}
\olsection{Teorema Wangun Normal}

Upamane kita bisa njlentrehake definisi fungsi
komputabel lan mriksa apa jlentrehan kang dianggep
cathetan lengkap komputasi fungsi iku ing sawenehing
input pancen bener. Mula, tanpa gumantung marang
model komputasi, kita mesthine bisa nemtokake nilai
fungsi komputabel~$f$ ing input~$x$ kanthi cara iki:
\begin{enumerate}
  \item Telusurana kabeh jlentrehan cathetan
  komputasi kang bisa ana.
  \item Priksanen apa cathetan tartamtu iku cathetan
  komputasi~$f(x)$.
  \item Yen wis ketemu cathetan kang bener, jupuk
  nilai~$f(x)$ saka cathetan iku.
\end{enumerate}
Kasunyatan yen cara iki pancen bener iku isi
teorema wangun normal Kleene.

\begin{thm}[Teorema Wangun Normal Kleene]
\ollabel{thm:normal-form}
Ana relasi rekursif primitif~$T(e, x, s)$ lan
fungsi rekursif primitif~$U(s)$, kanthi sipat iki:
yen $f$ fungsi komputabel parsial sembarang,
mula kanggo sawenehing~$e$,
\[
f(x) \simeq U(\umin{s}{T(e, x, s)})
\]
kanggo saben~$x$. Cathetan panyunting: tampilan iki kanggo fungsi unèr; fungsi mawa luwih saka siji argumen diwaca nganggo kode tuple input, kanthi pangodean kang cocog.
\end{thm}

\begin{proof}[Sketsa Bukti]
Kanggo model komputasi sembarang, jlentrehan
fungsi komputabel~$f$ bisa ditegesi kanthi tliti
lan dikode nganggo wilangan asli~$e$. Pangerten
bab ``rerangken komputasi'' kang nyathet proses
ngitung fungsi mawa indeks~$e$ ing input~$x$
uga bisa ditegesi kanthi tliti. Rerangken
komputasi kaya mangkono uga bisa dikode
minangka wilangan~$s$. Iki bisa ditindakake
kanthi cara kang ndadekake
\begin{enumerate}
  \item relasi $T(e, x, s)$, kang lumaku yen lan
  mung yen wilangan~$s$ ngode rerangken komputasi
  fungsi mawa indeks~$e$ ing input~$x$, lan
  \item fungsi $U(s)$ kang masangake rerangken
  komputasi mawa kode~$s$ karo asil pungkasan
  saka komputasi iku
\end{enumerate}
kalorone komputabel. Malah, relasi~$T$ lan
fungsi~$U$ iku rekursif primitif.
\end{proof}

\begin{explain}
Kanggo menehi bukti tliti tumrap Teorema Wangun
Normal, kita kudu netepake model komputasi,
nggarap pangodean jlentrehan fungsi komputabel
lan rerangken komputasi kanthi rinci, banjur
mriksa yen relasi~$T$ lan fungsi~$U$ iku
rekursif primitif. Kanggo akèh panganggone,
cukup yen $T$ lan~$U$ komputabel lan $U$~total.

Mbokmenawa bukti teorema wangun normal paling
gampang dielingi minangka ukara ringkes:
$\umin{s}{T(e, x, s)}$ nggoleki rerangken
komputasi fungsi mawa indeks~$e$ ing input~$x$,
lan $U$ ngasilake asil rerangken komputasi iku
yen rerangken kaya mangkono bisa ditemokake.
\end{explain}

Yen model komputasine yaiku fungsi rekursif
parsial (kang wiwitane digunakake Kleene),
teorema iki nuduhake yen mung siji panganggone
panggolekan tanpa wates, yaiku mung siji operator
$\umin{y}{f(\vec x, y)}$, kang dibutuhake
kanggo netepake fungsi sembarang. Ing pangerten
iki, saben fungsi rekursif parsial nduwé
wangun normal.

$T$ lan $U$ bisa digunakake kanggo netepake
enumerasi $\cfind{0}$, $\cfind{1}$,
$\cfind{2}$, \dots. Wiwit saiki, kita nganggep
wis netepake pilihan $T$ lan~$U$ kang trep,
lan njupuk persamaan
\[
\cfind{e}(x) \simeq U(\umin{s}{T(e,x,s)})
\]
minangka \emph{definisi} $\cfind{e}$.

Iki kasunyatan migunani liyane:

\begin{thm}
Saben fungsi komputabel parsial nduwé
indeks tanpa wates cacahé. Cathetan panyunting: ing kene digunakake pangindeksan program baku kang menehi kode beda kanggo instruksi beda lan ngidini instruksi tambahan tanpa ngowahi fungsi. Anane evaluator universal wae durung njamin sipat cacah indeks iki.
\end{thm}

Maneh, iki cetha sacara intuitif. Yen diwenehi
sawijining (jlentrehan) fungsi komputabel,
kita bisa nggawe jlentrehan liya kang ngitung
fungsi kang padha (pasangan input-asil padha),
nanging, umpamane, luwih dhisik nindakake
perkara kang ora mengaruhi komputasi (umpamane,
mriksa apa $0 = 0$, utawa ngetung nganti $5$,
lan sapiturute). Indeks jlentrehan kang diowahi
iku mesthi beda karo indeks asline. Kalorone
indeks saka fungsi kang padha, mung carane
ngitung rada beda.

\end{document}
